\section{Fundamental Theorem for normal matrix product states}%
\label{sec:peps_normal_mps}

The closed chain is the one-dimensional PEPS: the cycle graph on $n$ vertices
carrying one tensor of $D \times D$ matrices per site, so that a physical
configuration $\sigma = (i_1, \ldots, i_n)$ receives the amplitude
$\tr(A^{i_1} \cdots A^{i_n})$. The Fundamental Theorem for normal PEPS
therefore specializes to matrix product states once its blocking hypotheses are
read off the chain, and this section records the resulting statements,
following the matrix-product-state corollaries of
\cite[Section~4]{Molnar2018NormalPEPS} and the sharper system-size bound in its
closing part \emph{Normal MPS: alternative proof}.

The development proceeds in four steps. The closed chain is first identified
with a cycle of tensors, and the normal blocking hypotheses are obtained from
the injectivity of every arc of $L$ consecutive sites, giving the Fundamental
Theorem for normal matrix product states on a ring of $n \ge 3L$ sites with one
invertible gauge per bond. The translation-invariant case then collapses the
per-bond gauges to one gauge $Z$ and one scalar $\lambda$, with the length-$n$
trace identity forcing $\lambda^n = 1$.

For the sharper bound, the matrix-product-state form of
\cite[Lemma~5]{Molnar2018NormalPEPS} compares the $L + 1$ overlapping
windows around a fixed bond. If $C_j$ is the deformation of the $j$-th
window, then for length-$L$ words $a$ and $b$ the family of
consecutive-window trace identities strips to
$C_0(a)\,B^b = B^a\,C_L(b)$, and hence to
$C_0(a) = B^a X$ and $C_L(b) = X B^b$.
This is the step that lowers the system size from $n \ge 3L$ to
$n \ge 2L + 1$. Finally the same argument is run for site-dependent
chains, one tensor per site, which is the setting of
\cite[Lemma~5]{Molnar2018NormalPEPS}.

\subsection{Matrix product states on the closed chain}

\begin{theorem}[Cycle blocking hypotheses from arc injectivity]%
    \label{thm:peps_cycleNormalPEPSBlockingHypotheses}
    \lean{TNLean.PEPS.cycleArcFrom}
    \lean{TNLean.PEPS.NormalCycleArcInjectivityHypotheses}
    \lean{TNLean.PEPS.cycleEdgeBlockingData}
    \lean{TNLean.PEPS.cycleOneSiteSeparation}
    \lean{TNLean.PEPS.cycleNormalPEPSBlockingHypotheses}
    \lean{TNLean.PEPS.isCrossingEdge_cycleEdge_iff}
    \leanok
    \uses{def:peps_normalPEPSBlockingHypotheses,
        def:peps_regionInjectivityData}
    A closed chain of $n$ sites is the cycle graph on $n$ vertices, and a
    block of consecutive sites is an arc: the vertices reached clockwise
    from a starting site. Suppose $0 < L$, $3L \le n$, every arc
    of $L$ consecutive sites is injective for a region-injectivity
    assignment, and injectivity is closed under unions. Then the chain
    satisfies the normal blocking hypotheses: around every edge it blocks
    into the $L$ sites on one side, the $L$ sites on the other, and the
    remaining $n - 2L \ge L$ sites, all injective and covering the chain;
    and at every site the $L + 1$ sites starting there and the $L$ sites
    starting at its successor are injective regions with injective
    complements differing exactly at that site. Moreover the only
    nearest-neighbour pair joining the two blocks around an edge is the edge
    itself, since the far ends of the two blocks are separated by the
    remaining $n - 2L \ge L$ sites.
\end{theorem}
\begin{proof}\leanok
    Membership in an arc is a cyclic-distance inequality, so the
    disjointness, covering, and unique-crossing claims are modular
    arithmetic on the distances. Arcs of length between $L$ and $n$ are
    unions of arcs of length $L$, hence injective by union closure; this
    covers the two complements at a site and the third block around an
    edge.
\end{proof}

\begin{corollary}[Fundamental Theorem for normal MPS on a closed chain]%
    \label{cor:fundamentalTheorem_normalMPS_cycle}
    \lean{TNLean.PEPS.fundamentalTheorem_normalMPS_cycle}
    \leanok
    \uses{def:peps_regionInjectivityDataOf, def:GaugeEquivPEPS}
    Let $A$ and $B$ be tensors on the closed chain of $n$ sites (the cycle
    graph on $n$ vertices, one tensor per site with two virtual legs), with
    the same positive physical dimension and positive bond dimensions, and
    let $0 < L$ with $n \ge 3L$. Suppose blocking any $L$ consecutive sites
    gives an injective tensor, for $A$ and for $B$, and the two MPS generate
    the same state. Then the defining tensors are related by a local gauge:
    there are invertible matrices, one on each bond, whose endpoint actions
    transform $A_i$ into $B_i$ at every site with no leftover scalar. This
    is the source's family $Z_i$ ($i = 1, \ldots, n$, $n + 1 \equiv 1$) with
    $B_i = Z_i^{-1} A_i Z_{i+1}$, the first corollary following the general
    normal-PEPS theorem in \cite[Section~4]{Molnar2018NormalPEPS}, the
    statement for MPS without translation invariance at a fixed system size.
    The positive physical and bond dimensions are the counterexample-backed
    corrections carried over from the injective Fundamental Theorem
    (Theorem~\ref{thm:fundamentalTheorem_PEPS}), and $0 < L$ is implicit in
    the source's blocking of $L$ consecutive sites; the equality of the bond
    dimensions is not assumed but derived, and the connectivity of the cycle
    is proved, not assumed.
\end{corollary}
\begin{proof}\leanok
    \uses{thm:peps_cycleNormalPEPSBlockingHypotheses,
        thm:fundamentalTheorem_normalPEPS_connected,
        lem:peps_injectiveRegion_union_pos}
    Positive bond dimensions give union closure of injective regions for
    each tensor (Lemma~\ref{lem:peps_injectiveRegion_union_pos}), hence for
    the doubly injective assignment. The arc-injectivity hypothesis then
    yields the normal blocking hypotheses on the cycle with the unique
    crossing property
    (Theorem~\ref{thm:peps_cycleNormalPEPSBlockingHypotheses}), the cycle is
    connected, and the Fundamental Theorem for normal PEPS on a connected
    graph (Theorem~\ref{thm:fundamentalTheorem_normalPEPS_connected})
    supplies the scalar-free local gauge.
\end{proof}

\begin{corollary}[Uniqueness of the closed-chain gauges up to a constant]%
    \label{cor:fundamentalTheorem_normalMPS_cycle_gauge_unique}
    \lean{TNLean.PEPS.fundamentalTheorem_normalMPS_cycle_gauge_unique}
    \leanok
    \uses{def:peps_regionInjectivityDataOf, def:gaugeVertex}
    In the setting of
    Corollary~\ref{cor:fundamentalTheorem_normalMPS_cycle}, with the bond
    dimensions of the two tensors equal and positive and the arcs of $L$
    consecutive sites injective for each tensor, any two families of
    invertible bond matrices realizing the scalar-free gauge relation are
    proportional at every bond: $Z'_i = c \cdot Z_i$ for a nonzero scalar
    $c$ depending on the bond. This is the uniqueness clause of the first
    corollary following the general normal-PEPS theorem of
    \cite[Section~4]{Molnar2018NormalPEPS}: the gauges $Z_i$ are unique up
    to a multiplicative constant. Neither equality of the two states nor
    connectivity is needed, as in the general uniqueness clause.
\end{corollary}
\begin{proof}\leanok
    \uses{thm:peps_cycleNormalPEPSBlockingHypotheses,
        thm:fundamentalTheorem_normalPEPS_gauge_unique,
        lem:peps_injectiveRegion_union_pos}
    The arc-injectivity hypothesis and union closure from the positive bond
    dimensions (Lemma~\ref{lem:peps_injectiveRegion_union_pos}) assemble the
    normal blocking hypotheses on the cycle
    (Theorem~\ref{thm:peps_cycleNormalPEPSBlockingHypotheses}), and the
    per-edge uniqueness of the general theorem
    (Theorem~\ref{thm:fundamentalTheorem_normalPEPS_gauge_unique}) gives the
    proportionality at every bond.
\end{proof}

\begin{definition}[Cycle tensor of a matrix tensor]%
    \label{def:cycleTensorOfMPS}
    \lean{TNLean.PEPS.cycleSuccEdge}
    \lean{TNLean.PEPS.cycleLeftIncident}
    \lean{TNLean.PEPS.cycleRightIncident}
    \lean{TNLean.PEPS.cycleTensorOfMPS}
    \leanok
    \uses{def:peps_tensor, def:mps_tensor}
    Every edge of the cycle graph on $n \ge 3$ vertices joins a site to its
    cyclic successor, so sites and bonds are in bijection, and every site has
    exactly two incident bonds: the bond shared with its cyclic predecessor
    and the bond shared with its cyclic successor. The \emph{cycle tensor}
    of a tensor $A = (A^i)_{i=0}^{d-1}$ of $D \times D$ matrices places one
    copy of $A$ at every site of the closed chain of $n$ sites, every bond of
    dimension $D$, with the row index of $A^i$ read from the bond shared with
    the predecessor and the column index from the bond shared with the
    successor.
\end{definition}

\begin{lemma}[State coefficients of the cycle tensor]%
    \label{lem:stateCoeff_cycleTensorOfMPS}
    \lean{MPSTensor.evalWord_ofFn_apply}
    \lean{MPSTensor.trace_evalWord_eq_sum_cyclic}
    \lean{TNLean.PEPS.stateCoeff_cycleTensorOfMPS}
    \leanok
    \uses{def:cycleTensorOfMPS, def:peps_stateCoeff, def:mpv}
    For $n \ge 3$ and every configuration
    $\sigma = (i_1, \ldots, i_n)$, the state coefficient of the cycle
    tensor of $A$ at $\sigma$ is the matrix product vector component
    $V^{(n)}(A)_\sigma = \tr\!(A^{i_1} \cdots A^{i_n})$.
\end{lemma}
\begin{proof}\leanok
    An entry of a word product is a sum over paths of bond indices pinned at
    the two ends of products of matrix entries read along the path, by
    induction on the word; closing the two pinned ends into a loop, the trace
    of a closed word product is the sum over all cyclic bond configurations
    of the products of entries around the loop. Contracting the cycle tensor
    sums the per-site entries over all assignments of bond indices, which is
    that cyclic sum after matching each bond with the site on its
    predecessor side.
\end{proof}

\begin{lemma}[Arc injectivity of the cycle tensor from block injectivity]%
    \label{lem:regionBlockedTensorInjective_cycleTensorOfMPS}
    \lean{TNLean.PEPS.isRegionBoundaryEdge_cycleArcFrom_iff}
    \lean{TNLean.PEPS.regionBlockedWeight_cycleTensorOfMPS}
    \lean{TNLean.PEPS.matrix_eq_zero_of_isNBlkInjective}
    \lean{TNLean.PEPS.regionBlockedTensorInjective_cycleTensorOfMPS}
    \leanok
    \uses{def:cycleTensorOfMPS, def:l_blk_injective,
        def:peps_regionInjectivityDataOf}
    Let $0 < L < n$ with $n \ge 3$, let the bond dimension be positive, and
    suppose $A$ is $L$-block injective. Then for every arc of $L$
    consecutive sites of the closed chain, the blocked tensor of the arc of
    the cycle tensor of $A$ is injective.
\end{lemma}
\begin{proof}\leanok
    An arc of $L < n$ consecutive sites has exactly two boundary-crossing
    bonds, the bond entering its first site and the bond leaving its last
    site, by cyclic-distance arithmetic. The blocked weight of the arc at a
    boundary assignment $(a, b)$ and a physical word $x$ is
    $D^{\,n - (L+1)}\,(A^{x_1} \cdots A^{x_L})_{a b}$: the bonds inside
    the arc contract to the matrix product, the two boundary bonds stay
    pinned, and each bond away from the arc contributes a free factor $D$.
    A linear relation among the blocked components over the boundary pairs
    is a matrix $C$ with $\tr\!(C^{\mathsf T} A^{x_1} \cdots A^{x_L}) = 0$
    for every word $x$; the word products span the full matrix algebra by
    $L$-block injectivity, and the trace pairing is nondegenerate, so
    $C = 0$.
\end{proof}

\subsection{Translation-invariant matrix product states}

When the same tensor sits at every site, the closed chain ties the per-bond
gauges together. If one gauge $Z$ and one scalar $\lambda$ satisfy
$B^i = \lambda\,Z^{-1}A^iZ$ for every physical index $i$, then for a
configuration $\sigma=(i_1,\ldots,i_n)$,
\begin{align}
    \tr(B^{i_1}\cdots B^{i_n})
    &= \lambda^n\,\tr(Z^{-1}A^{i_1}\cdots A^{i_n}Z)
     = \lambda^n\,\tr(A^{i_1}\cdots A^{i_n}).
    \notag
\end{align}
Equality of the matrix product states therefore forces $\lambda^n = 1$.
Thus $Z$ and $\lambda$ capture precisely the freedom in the
translation-invariant corollary of \cite[Section~4]{Molnar2018NormalPEPS}.

\begin{corollary}[Fundamental Theorem for translation-invariant normal MPS
        on a closed chain, matrix form]%
    \label{cor:fundamentalTheorem_normalMPS}
    \lean{TNLean.PEPS.fundamentalTheorem_normalMPS}
    \leanok
    \uses{def:l_blk_injective, def:mpv}
    Let $n$ be positive, let $0 < L$ with $2L + 1 \le n$, and let $A$ and $B$
    be tensors of $D \times D$ matrices with $D>0$, over physical dimension
    $d$, each
    $L$-block injective, whose matrix product vectors agree at system size
    $n$: $V^{(n)}(A)_\sigma = V^{(n)}(B)_\sigma$ for every configuration
    $\sigma$. Then there are invertible matrices $Z_v$, one per bond of the
    closed chain ($v = 1, \ldots, n$, $n + 1 \equiv 1$), with
    $B^i = Z_v^{-1}\, A^i\, Z_{v+1}$ for every site $v$ and every $i$.
    This is the first corollary following the general normal-PEPS theorem of
    \cite[Section~4]{Molnar2018NormalPEPS} specialized to one
    site-independent tensor: the source states the corollary for
    site-dependent families, and its conclusion is the per-bond gauge family
    delivered here. The system size is the optimal $n \ge 2L + 1$ of the
    source's alternative proof (its Section ``Normal MPS: alternative
    proof''), rather than the $n \ge 3L$ of the blocking route. The source's
    translation-invariant corollary (a single gauge $Z$ and a constant
    $\lambda$ with $\lambda^n = 1$) refines this conclusion and is not part
    of this statement. A vanishing physical dimension is already excluded
    by block injectivity.
\end{corollary}
\begin{proof}\leanok
    \uses{cor:fundamentalTheorem_normalMPS_of_overlap}
    The per-bond form of the strengthened closed-chain corollary
    (Corollary~\ref{cor:fundamentalTheorem_normalMPS_of_overlap}), proved
    through the overlapping windows at $n \ge 2L + 1$, gives the per-bond
    gauge family directly.
\end{proof}

\begin{corollary}[Uniqueness of the per-bond gauges up to one constant,
        matrix form]%
    \label{cor:fundamentalTheorem_normalMPS_gauge_unique}
    \lean{TNLean.PEPS.fundamentalTheorem_normalMPS_gauge_unique}
    \leanok
    \uses{def:l_blk_injective}
    Let $n$ be positive, let $0 < L$ with $3L \le n$, and let $A$ and $B$ be
    tensors of $D \times D$ matrices with $D>0$, each $L$-block injective.
    If two
    families $Z$ and $Z'$ of invertible matrices on the bonds of the closed
    chain both satisfy $B^i = Z_v^{-1} A^i Z_{v+1}$ at every site, then they
    are proportional by a single constant: there is a nonzero scalar $c$
    with $Z'_v = c\,Z_v$ at every bond. This is the uniqueness clause of
    the first corollary following the general normal-PEPS theorem of
    \cite[Section~4]{Molnar2018NormalPEPS}, for one site-independent
    tensor. Equality of the two states is not needed.
\end{corollary}
\begin{proof}\leanok
    \uses{cor:fundamentalTheorem_normalMPS_cycle_gauge_unique,
        lem:regionBlockedTensorInjective_cycleTensorOfMPS,
        def:cycleTensorOfMPS, def:gaugeVertex}
    Each per-bond family converts back into a per-edge family - the
    transpose of the bond gauge away from the seam, the inverse on the seam
    edge - realizing the per-site gauge identity between the cycle tensors,
    so the closed-chain uniqueness clause
    (Corollary~\ref{cor:fundamentalTheorem_normalMPS_cycle_gauge_unique})
    gives a constant on every edge; undoing the conversion gives a constant
    per bond. The relation $B^i = Z_v^{-1} A^i Z_{v+1}$ for both families
    forces the scalar $\mu_v^{-1}\mu_{v+1}$ to fix the nonzero tensor $B$,
    so consecutive constants agree, and going around the cycle merges them
    into one.
\end{proof}

\begin{lemma}[Consecutive per-bond gauges are proportional]%
    \label{lem:gaugeFamily_succ_proportional}
    \lean{TNLean.PEPS.gaugeFamily_succ_proportional}
    \leanok
    \uses{def:l_blk_injective, def:eval_word}
    Let $n$ be positive, let $A$ and $B$ be tensors of $D \times D$ matrices
    over physical dimension $d$ with $A$ being $L$-block injective, and let
    $Z_v$ ($v = 1, \ldots, n$, $n + 1 \equiv 1$) be invertible matrices with
    $B^i = Z_v^{-1} A^i Z_{v+1}$ at every site $v$ and every $i$. Then for
    every site $v$ there is a nonzero scalar $c$ with $Z_{v+1} = c\,Z_v$.
\end{lemma}
\begin{proof}\leanok
    \uses{def:eval_word}
    Iterating the per-bond relation along a word $x$ of length $L$ from the
    starting site $v$ gives
    $B^{x_1} \cdots B^{x_L} = Z_v^{-1} A^{x_1} \cdots A^{x_L} Z_{v+L}$.
    Comparing the iterations from $v$ and from $v + 1$ over the word
    products $A^{x_1} \cdots A^{x_L}$, which span the full matrix algebra by
    $L$-block injectivity, the comparison extends linearly to every matrix
    $W$: $Z_v^{-1} W Z_{v+L} = Z_{v+1}^{-1} W Z_{v+1+L}$. The empty word
    ($W$ the identity) pins the two bond transports
    $Z_v^{-1} Z_{v+L} = Z_{v+1}^{-1} Z_{v+1+L}$ to the same invertible
    matrix; cancelling it leaves $Z_v^{-1} W Z_v = Z_{v+1}^{-1} W Z_{v+1}$
    for every $W$, so $Z_{v+1} Z_v^{-1}$ commutes with the full matrix
    algebra, hence is a scalar, nonzero by invertibility.
\end{proof}

\begin{corollary}[Fundamental Theorem for translation-invariant normal MPS
        on a closed chain, single-gauge form]%
    \label{cor:fundamentalTheorem_normalMPS_translationInvariant}
    \lean{TNLean.PEPS.fundamentalTheorem_normalMPS_translationInvariant}
    \leanok
    \uses{def:l_blk_injective, def:mpv}
    Let $n$ be positive, let $0 < L$ with $2L + 1 \le n$, and let $A$ and $B$
    be tensors of $D \times D$ matrices with $D>0$, over physical dimension
    $d$, each
    $L$-block injective, whose matrix product vectors agree at system size
    $n$: $V^{(n)}(A)_\sigma = V^{(n)}(B)_\sigma$ for every configuration
    $\sigma$. Then there are an invertible matrix $Z$ and a constant
    $\lambda$ with $\lambda^n = 1$ such that
    $B^i = \lambda\, Z^{-1} A^i Z$ for every $i$.
    This is the translation-invariant corollary of
    \cite[Section~4]{Molnar2018NormalPEPS}, the statement for TI MPS
    following the first corollary after the general normal-PEPS theorem. The
    system size is the optimal $n \ge 2L + 1$ of the source's alternative
    proof.
\end{corollary}
\begin{proof}\leanok
    \uses{cor:fundamentalTheorem_normalMPS_translationInvariant_of_overlap}
    This is the single-gauge form of the strengthened closed-chain corollary
    (Corollary~\ref{cor:fundamentalTheorem_normalMPS_translationInvariant_of_overlap}),
    proved through the overlapping windows at $n \ge 2L + 1$.
\end{proof}

\begin{corollary}[Uniqueness of the translation-invariant gauge up to a
        constant]%
    \label{cor:fundamentalTheorem_normalMPS_translationInvariant_gauge_unique}
    \lean{TNLean.PEPS.fundamentalTheorem_normalMPS_translationInvariant_gauge_unique}
    \leanok
    \uses{def:l_blk_injective}
    Let $0 < L$ and $0<D$, and let $A$ and $B$ be tensors of $D \times D$ matrices,
    each $L$-block injective. If invertible matrices $Z$, $Z'$ and
    constants $\lambda$, $\lambda'$ satisfy
    $B^i = \lambda\, Z^{-1} A^i Z$ and $B^i = \lambda'\, Z'^{-1} A^i Z'$ for
    every $i$, then there is a nonzero scalar $c$ with $Z' = c\,Z$. This is
    the uniqueness clause of the translation-invariant corollary of
    \cite[Section~4]{Molnar2018NormalPEPS} for TI MPS. No system size
    enters the statement: equality of the two states is not needed, nor are
    the root-of-unity conditions on $\lambda$ and $\lambda'$.
\end{corollary}
\begin{proof}\leanok
    \uses{def:eval_word}
    The scalar $\lambda$ is nonzero because $B$ is $L$-block injective with
    positive bond dimension, so some $B^{i_0} \neq 0$. Iterating each relation
    along a word $x$ of length $L$ gives
    $B^{x_1} \cdots B^{x_L} = \lambda^L\, Z^{-1} A^{x_1} \cdots A^{x_L} Z$
    and the primed analog; the word products span the full matrix algebra
    by $L$-block injectivity, so
    $\lambda^L\, Z^{-1} W Z = \lambda'^L\, Z'^{-1} W Z'$ for every matrix
    $W$. The identity pins $\lambda^L = \lambda'^L$; cancelling the
    nonzero factor leaves $Z^{-1} W Z = Z'^{-1} W Z'$ for every $W$, so
    $Z' Z^{-1}$ commutes with the full matrix algebra, hence is a nonzero
    scalar.
\end{proof}

\subsection{\texorpdfstring{The optimal system size $n \ge 2L + 1$}{The optimal system size n >= 2L+1}}

The bound $n \ge 3L$ is not optimal. A bond of the closed chain is straddled by
the $L + 1$ overlapping windows of $L$ sites around it. For consecutive
windows, the deformed closed-chain coefficients have the form
\begin{align}
    \tr(B^{p}\, C_j(u_1 \cdots u_L)\, B^{u_{L+1}}\, B^{s})
    &=
    \tr(B^{p}\, B^{u_1}\, C_{j+1}(u_2 \cdots u_{L+1})\, B^{s}).
    \notag
\end{align}
Injectivity on the complementary arc strips these identities to
\begin{align}
    C_j(u_1 \cdots u_L)\,B^{u_{L+1}}
    &=
    B^{u_1}\,C_{j+1}(u_2 \cdots u_{L+1}),
    \notag
\end{align}
and chaining the relations across the $L + 1$ windows gives
$C_0(a)\,B^b = B^a\,C_L(b)$.
This is the overlap mechanism in \cite[Lemma~5]{Molnar2018NormalPEPS}.
The lemmas below assemble that argument: block injectivity extends to every
length, word products transport between equal-state chains, and the bond
operator is extracted and shown to act by conjugation, lowering the system size
to $n \ge 2L + 1$.

\begin{lemma}[Block injectivity extends to longer blocks]%
    \label{lem:isNBlkInjective_of_le}
    \lean{MPSTensor.isNBlkInjective_of_le}
    \leanok
    \uses{def:l_blk_injective}
    Let $A$ be a tensor of $D \times D$ matrices over physical dimension
    $d$, $L$-block injective with $0 < L$. Then $A$ is $m$-block injective
    for every $m \ge L$. This is the chain form of the statement that any
    region of at least the blocking size is injective, derived in the
    closing section of \cite{Molnar2018NormalPEPS} (``Normal MPS:
    alternative proof'') from its union lemma.
\end{lemma}
\begin{proof}\leanok
    \uses{def:eval_word}
    It suffices to pass from length $m$ to $m + 1$. Write the identity as
    a combination of length-$L$ word products and peel the leading letter
    of each: any matrix $P$ becomes a combination of terms
    $A^{i}\,(W P)$ with $W$ a word product of length $L - 1$. Each factor
    $W P$ is a combination of length-$m$ word products by the induction
    hypothesis, so each term is a combination of length-$(m+1)$ word
    products.
\end{proof}

\begin{lemma}[Word transport between same-state networks]%
    \label{lem:exists_mpvTransport}
    \lean{MPSTensor.exists_mpvTransport}
    \lean{MPSTensor.exists_evalWordTransport}
    \leanok
    \uses{def:l_blk_injective, def:mpv, def:eval_word}
    Let $A$ and $B$ be tensors of $D \times D$ matrices over physical
    dimension $d$, let $n = k + q$, and suppose the length-$k$ and
    length-$q$ word products of $A$ each span the full matrix algebra. If
    the matrix product vectors of $A$ and $B$ agree at system size $n$,
    then there is a linear map $\Lambda$ of the matrix algebra with
    $\Lambda(A^{x_1} \cdots A^{x_k}) = B^{x_1} \cdots B^{x_k}$ for every
    word $x$ of length $k$. The same holds when, instead of the traces,
    the word products of the two tensors are themselves equal at length
    $n$. This is the operator-transport mechanism underlying Lemma~5 of
    \cite{Molnar2018NormalPEPS}: an operator expanded over the window
    products of one network is re-read in the other.
\end{lemma}
\begin{proof}\leanok
    \uses{def:mpv, def:eval_word}
    Pair the matrix algebra against the complementary word products: let
    $\Psi_A(M)$ be the family of traces
    $\tr(M\,A^{y_1} \cdots A^{y_q})$ over the length-$q$ words $y$, and
    $\Psi_B$ its $B$-analog. $\Psi_A$ is injective because the length-$q$
    products span and the trace pairing is nondegenerate. On the spanning
    family of length-$k$ products the two pairings match:
    $\Psi_A(A^{x}) = \Psi_B(B^{x})$, since both compute closed-chain
    coefficients of total length $n$, equal by hypothesis. Hence the range
    of $\Psi_A$ lies in that of $\Psi_B$, forcing $\Psi_B$ to be injective
    by dimension count, and a left inverse of $\Psi_B$ composed with
    $\Psi_A$ is the transport. For the matrix-product version replace the
    traces by the products themselves.
\end{proof}

\begin{lemma}[Bond-operator extraction from overlapping windows]%
    \label{lem:overlapWindow_exists_bondOperator}
    \lean{MPSTensor.overlapWindow_exists_bondOperator}
    \lean{MPSTensor.bondOperator_unique}
    \leanok
    \uses{def:l_blk_injective, def:eval_word}
    Let $0 < L$ with $2L + 1 \le n$, let $B$ be a tensor of $D \times D$
    matrices over physical dimension $d$, $L$-block injective, and let
    $C_0, \ldots, C_L$ be deformed window tensors: each $C_j$ assigns a
    matrix to every word of length $L$, read as replacing the blocked
    tensor of the $L$-site window starting $j$ sites left of a fixed bond
    of the closed chain on $n$ sites. Suppose the deformed states agree
    for consecutive window positions: for every $j < L$, every prefix $p$
    of $j$ sites, every assignment $u$ on the $L + 1$ sites covered by the
    two windows, and every suffix $s$ on the remaining sites,
    \begin{align}
        \tr(B^{p}\, C_j(u_1 \cdots u_L)\, B^{u_{L+1}}\, B^{s})
        &=
        \tr\bigl(B^{p}\, B^{u_1}\, C_{j+1}(u_2 \cdots u_{L+1})
        B^{s}\bigr).
        \notag
    \end{align}
    Then the deformation is a virtual operation on the bond: there is a
    matrix $X$ such that, for all length-$L$ words $a$ and $b$,
    \begin{align}
        C_0(a) &= B^{a_1} \cdots B^{a_L}\, X,
        \notag\\
        C_L(b) &= X\, B^{b_1} \cdots B^{b_L},
        \notag
    \end{align}
    and $X$ is unique. This is Lemma~5
    of \cite{Molnar2018NormalPEPS} (its closing section, ``Normal MPS:
    alternative proof''), specialized to one site-independent tensor; the
    source states it for site-dependent tensors on five sites with $L = 2$,
    with the deformed window tensors arising as physical operators
    contracted with the blocked windows they act on; the site-dependent
    form, in the setting in which the source states the lemma, is
    Lemma~\ref{lem:chain_overlapWindow_exists_bondOperator}, and this
    statement is its constant specialization. The
    algebra-homomorphism clause of the source lemma is established with the
    insertion correspondence in
    Lemma~\ref{lem:exists_conjugation_of_mpv_eq}.
\end{lemma}
\begin{proof}\leanok
    \uses{lem:chain_overlapWindow_exists_bondOperator, def:eval_word}
    Place the tensor at every site of the site-dependent form
    (Lemma~\ref{lem:chain_overlapWindow_exists_bondOperator}): the arc
    products of the constant chain are the word products of the repeated
    tensor, and $L$-block injectivity of the tensor is window injectivity
    of the constant chain.
\end{proof}

\begin{lemma}[Conjugate word products at one length from Lemma~5]%
    \label{lem:exists_conjugation_of_mpv_eq}
    \lean{MPSTensor.exists_conjugation_of_mpv_eq}
    \leanok
    \uses{def:l_blk_injective, def:mpv, def:eval_word}
    Let $0 < L$ with $2L + 1 \le n$ and $0<D$, and let $A$ and $B$ be tensors of
    $D \times D$ matrices over physical dimension $d$, each $L$-block
    injective, whose matrix product vectors agree at system size $n$.
    Then there is an invertible matrix $Z$ such that, for every word $w$ of
    length $n$,
    \begin{align}
        B^{w_1} \cdots B^{w_n}
        &= Z\, A^{w_1} \cdots A^{w_n}\, Z^{-1}.
        \notag
    \end{align}
    This captures the insertion
    correspondence of Lemma~5 of \cite{Molnar2018NormalPEPS} and its
    algebra-homomorphism clause for one site-independent tensor: the map
    sending an insertion on one network to the corresponding insertion on
    the other is an automorphism of the matrix algebra, hence inner.
\end{lemma}
\begin{proof}\leanok
    \uses{cor:fundamentalTheorem_normalMPSChain_of_overlap,
        thm:mps_chain_gaugeEquiv_sameState}
    Regard $A$ and $B$ as constant site-dependent chains on $n$ sites.
    Window injectivity of a constant chain is block injectivity of the
    single tensor, and agreement of the matrix product vectors at size $n$
    is the same-state hypothesis, so the site-dependent overlapping-window
    capstone
    (Corollary~\ref{cor:fundamentalTheorem_normalMPSChain_of_overlap})
    supplies a per-bond gauge family relating the two chains. The arc
    products of a constant chain are the word products of the tensor, and
    the gauge conjugation of arc products
    (Theorem~\ref{thm:mps_chain_gaugeEquiv_sameState}), read at the full
    chain length $n$ where the end bond wraps to the starting bond, is
    conjugation of the length-$n$ word products by the gauge at the
    starting bond.
\end{proof}

\begin{lemma}[Rigidity of word products at one length]%
    \label{lem:proportional_of_evalWord_eq}
    \lean{MPSTensor.proportional_of_evalWord_eq}
    \leanok
    \uses{def:l_blk_injective, def:eval_word}
    Let $0 < L$ with $2L + 1 \le n$ and $0 < D$, and let $A$ and $C$ be
    tensors of $D \times D$ matrices over physical dimension $d$, each
    $L$-block injective, with equal word products at length $n$:
    $C^{w_1} \cdots C^{w_n} = A^{w_1} \cdots A^{w_n}$ for every word $w$ of
    length $n$. Then there is a constant $\lambda$ with $\lambda^n = 1$
    and $C^i = \lambda A^i$ for every $i$.
\end{lemma}
\begin{proof}\leanok
    \uses{lem:exists_mpvTransport, lem:isNBlkInjective_of_le,
        def:eval_word}
    Transport word products of $A$ to those of $C$ at the window length
    $L$, at $L + 1$, at the complementary length $n - 1 - L$, and at
    $n - L$ (Lemma~\ref{lem:exists_mpvTransport}, matrix-product version;
    the lengths lie between $L$ and $n - L$ by $2L + 1 \le n$). For a
    transport pair at lengths $k$ and $k + 1$, the transported identity
    $\Lambda_k(\Id)$ commutes with every letter of $C$ - its two one-letter
    extensions agree, $C^i\,\Lambda_k(\Id) = \Lambda_{k+1}(A^i) =
        \Lambda_k(\Id)\,C^i$ - hence with the full matrix algebra by block
    injectivity, so it is a scalar $c_k \Id$, nonzero because the transport
    is surjective onto the spanning word products, hence bijective.
    Peeling one letter off the full-length equality through the spanning
    windows of lengths $L$ and $n - 1 - L$ then leaves
    $C^i = (c_L\,c_{n-1-L})^{-1} A^i =: \lambda A^i$; evaluating any
    nonzero length-$n$ word product gives $\lambda^n = 1$.
\end{proof}

\begin{corollary}[Fundamental Theorem for translation-invariant normal MPS
        at $n \ge 2L + 1$, single-gauge form]%
    \label{cor:fundamentalTheorem_normalMPS_translationInvariant_of_overlap}
    \lean{TNLean.PEPS.fundamentalTheorem_normalMPS_translationInvariant_of_overlap}
    \leanok
    \uses{def:l_blk_injective, def:mpv}
    Let $0 < L$ with $2L + 1 \le n$ and $0<D$, and let $A$ and $B$ be tensors of
    $D \times D$ matrices over physical dimension $d$, each $L$-block
    injective, whose matrix product vectors agree at system size $n$:
    $V^{(n)}(A)_\sigma = V^{(n)}(B)_\sigma$ for every configuration
    $\sigma$. Then there are an invertible matrix $Z$ and a constant
    $\lambda$ with $\lambda^n = 1$ such that
    $B^i = \lambda\, Z^{-1} A^i Z$ for every $i$.
    This is the strengthened closed-chain corollary of
    \cite{Molnar2018NormalPEPS} (the corollary after Lemma~5 in its closing
    section ``Normal MPS: alternative proof''), proved through the
    overlapping windows at the optimal $n \ge 2L + 1$ rather than the
    $n \ge 3L$ of the blocking route; the single-gauge corollary
    \ref{cor:fundamentalTheorem_normalMPS_translationInvariant} delegates to
    it. The uniqueness clause of the source
    corollary - the gauge $Z$ is unique up to a multiplicative constant -
    needs no system size and is
    Corollary~\ref{cor:fundamentalTheorem_normalMPS_translationInvariant_gauge_unique}.
\end{corollary}
\begin{proof}\leanok
    \uses{lem:exists_conjugation_of_mpv_eq, lem:proportional_of_evalWord_eq}
    By Lemma~\ref{lem:exists_conjugation_of_mpv_eq} the word products of
    $B$ at length $n$ are those of $A$ conjugated by an invertible matrix.
    Absorbing the conjugation into $B$ - which preserves $L$-block
    injectivity - leaves two tensors with equal word products at length
    $n$, and Lemma~\ref{lem:proportional_of_evalWord_eq} makes them
    proportional letterwise by an $n$-th root of unity. Undoing the
    absorption gives $B^i = \lambda\, Z^{-1} A^i Z$.
\end{proof}

\begin{corollary}[Fundamental Theorem for translation-invariant normal MPS
        at $n \ge 2L + 1$, per-bond form]%
    \label{cor:fundamentalTheorem_normalMPS_of_overlap}
    \lean{TNLean.PEPS.fundamentalTheorem_normalMPS_of_overlap}
    \leanok
    \uses{def:l_blk_injective, def:mpv}
    Let $n$ be positive, let $0 < L$ with $2L + 1 \le n$ and $0<D$, and let
    $A$ and $B$ be tensors of $D \times D$ matrices over physical dimension $d$,
    each $L$-block injective, whose matrix product vectors agree at system
    size $n$. Then there are invertible matrices $Z_v$, one per bond of
    the closed chain ($v = 1, \ldots, n$, $n + 1 \equiv 1$), with
    $B^i = Z_v^{-1}\, A^i\, Z_{v+1}$ for every site $v$ and every $i$.
    This is the per-bond gauge family of the closed-chain corollary, proved
    at the source's strengthened bound $n \ge 2L + 1$
    (\cite{Molnar2018NormalPEPS}, the corollary after Lemma~5 in its
    closing section ``Normal MPS: alternative proof''). The named matrix-level
    Corollary~\ref{cor:fundamentalTheorem_normalMPS} delegates to it.
\end{corollary}
\begin{proof}\leanok
    \uses{cor:fundamentalTheorem_normalMPS_translationInvariant_of_overlap}
    The single-gauge form
    (Corollary~\ref{cor:fundamentalTheorem_normalMPS_translationInvariant_of_overlap})
    gives $Z$ and $\lambda$ with $\lambda^n = 1$ and
    $B^i = \lambda\,Z^{-1} A^i Z$. The family $Z_v := \lambda^{v} Z$
    dresses the gauge with powers of the root of unity: away from the seam
    the consecutive powers contribute the factor $\lambda$, and across the
    seam the wraparound contributes $\lambda^{1-n} = \lambda$ because
    $\lambda^n = 1$, so $B^i = Z_v^{-1} A^i Z_{v+1}$ at every site.
\end{proof}
